\documentclass[10pt,a4paper]{amsart}
\usepackage[margin=19mm]{geometry}
\usepackage{amsmath,amssymb,mathtools}
\usepackage[scr=rsfs,cal=euler]{mathalfa}
\usepackage{xfrac}
\usepackage[unicode,hidelinks,bookmarks=false]{hyperref}
\newcommand{\ie}{i.e.\ }
\newcommand{\cf}{cf.\ }
\newcommand{\resp}{respectively\ }
\newcommand{\Subsection}{\subsection}
\providecommand{\nobreakdash}{\nobreak\hbox{-}}
\setlength{\emergencystretch}{2em}
\multlinegap0pt
\usepackage{amssymb}
\usepackage{bm}
\usepackage{enumerate}
\usepackage{tikz}
\usepackage[shortlabels]{enumitem}
\usepackage[normalem]{ulem}
\newtheorem{lem}{Lemma}
\newtheorem{thm}{Theorem}
\DeclareMathOperator{\End}{End}
\DeclareMathOperator{\Td}{Td}
\def\a{\alpha}
\def\bl{\bar l}
\def\bp{\bar\partial}
\def\br{\bar r}
\DeclareMathOperator{\cc}{c}
\def\p{\partial}
\title{Partition functions of determinantal point processes on polarized K\"ahler manifolds}
\author{Kiyoon Eum}
\date{Source: arXiv:2503.01524v5 (CC BY 4.0)}
\begin{document}
\maketitle

\noindent\emph{Source and licence.} Partition functions of determinantal point processes on polarized K\"ahler manifolds. Version arXiv:2503.01524v5 (CC BY 4.0), distributed by arXiv under the Creative Commons Attribution 4.0 International licence, http://creativecommons.org/licenses/by/4.0/, as recorded in the arXiv licence metadata for this exact version and shown on the abstract page https://arxiv.org/abs/2503.01524v5. Full licence notice: \texttt{licenses/arxiv-2503.01524-CC-BY-4.0.txt}.\par\smallskip

\noindent\textit{Editorial scope.} Counted unit: the article's thm futbar (source0.tex lines 706--716). The article's complete original proof of it is delivered with it (source0.tex lines 718--795, one \texttt{proof} environment of 78 lines). No mathematics is written, repaired or re-derived: the statement and the proof are the article's own lines, in the article's own notation, and no retained line is substituted or re-flowed.  Retained uncounted as the source-local context the proof needs: lines 676--686.  Nothing was omitted from any retained span, so the package carries no omission record.  No retained line contains a citation, so the wrapper carries no bibliography and no bibliography entry.  No retained line contains a figure environment, an image-inclusion command or drawing code, so no image is delivered and the package declares no asset dependency.  Editorial formatting only, in addition: an AMS \texttt{amsart} preamble that restates the article's own declarations and macros used by the retained lines, namely the environments \texttt{lem}, \texttt{thm} on separate counters, together with the standard packages those lines need.  The retained environments print in this document as Lemma 1, Theorem 1 in order of appearance, whereas the article prints the same items with the numbering of its own full document.  The complete native source of the article as distributed in the arXiv e-print for this exact version is bundled unchanged as \texttt{sources/174144/source0.tex}.
\begin{lem}\label{lulemm}
We have
\begin{equation}\label{nabx}
\iota_X R = -\bp \nabla X
\end{equation}
and 
\begin{equation}\label{nabx2}
\iota_{\overline{X}} \overline{R} = \p \overline{\nabla X}.
\end{equation}

\end{lem}

\begin{thm}\label{futbar}
Under the same notations as above, for $j\geq 0$, 
\begin{equation}\label{FX}
\widetilde{F_j}(\omega,X):=\int_M \Td_j(R+\nabla X)(\omega+\theta_X)^{n+1-j} 
\end{equation}
and
\begin{equation}\label{FbarX}
\widetilde{F_j}(\omega,\overline{X}):=\int_M \Td_j(R+(\overline{\nabla X})^{\flat,\sharp})(\omega-\overline{\theta_X})^{n+1-j} 
\end{equation}
are independent of the choice of the K\"ahler metric in $\cc_1(L)$, under the normalization of $\theta_X$ by $\int_M \theta_X \omega^n =0$. 
\end{thm}

\begin{proof}
We first prove the claim for (\ref{FX}). The proof for (\ref{FbarX}) is exactly the same, as we will explain.

Let $\omega_t=\omega+i\p\bp \varphi_t$ be an arbitrary one-parameter family of K\"ahler forms in $[\omega]$ with $\omega_0=\omega$. We will show that $\frac{d}{dt}\widetilde{F_j}(\omega_t,X)=0$. Let $\theta_t$ be a holomorphy potential of $X$ with respect to $\omega_t$, that is, $\iota_X \omega_t = -\bp \theta_t$. Setting $\a_t:=-i\p\dot{\varphi_t}$, we have $\dot{\omega_t}=i\p\bp\dot{\varphi_t}=\bp \a_t$.
Since 
\begin{equation*}
\iota_X \omega_t=\iota_X \omega + i\iota_X\p\bp\varphi_t=-\bp\theta_0+i\bp\iota_X (\p\varphi_t), 
\end{equation*}
we have $\theta_t=\theta_0-i\iota_X \p\varphi_t$ and thus $\dot{\theta_t}=\iota_X \a_t$ up to an additive constant.
We used the fact that $X\in\mathfrak{h}(M)$. Also, by torsion freeness of the Levi–Civita connection we have
\begin{equation*}
\dot{\left(\nabla X \right)}=\dot{X^q_p}=\dot{\Gamma^{q}_{pr}}X^r= \iota_X \dot{\Gamma_t},
\end{equation*}
where $\Gamma_t$ denotes the Levi–Civita connection $1$-form for $\omega_t$. Note that $\dot{\Gamma_t}$ is globally well defined and $\dot{R_t}=\bp\dot{\Gamma_t}$. Using these, we compute $\frac{d}{dt}\widetilde{F_j}(\omega_t,X)$ (we suppress the subscript $t$):  
\begin{align}\label{5}
&\frac{d}{dt}\widetilde{F_j}(\omega_t,X)=\int_M j\Td_j (\dot{\nabla X}+\dot{R}, \nabla X + R,\cdots)(\omega+\theta)^{n+1-j}\nonumber\\
&\qquad\qquad\qquad\quad+(n+1-j)\Td_j(\nabla X + R)(\dot{\theta}+\dot{\omega})(\omega+\theta)^{n-j} \nonumber\\
&=\int_M j\Td_j (\iota_X \dot{\Gamma}+\bp\dot{\Gamma}, \nabla X + R,\cdots)(\omega+\theta)^{n+1-j} \nonumber \\
&\qquad +(n+1-j)\Td_j(\nabla X + R)(\iota_X \a+\bp \a)(\omega+\theta)^{n-j}\nonumber\\
&=\int_M j\Td_j (\iota_X \dot{\Gamma}, \nabla X + R,\cdots)(\omega+\theta)^{n+1-j} \nonumber\\
&\qquad+j(j-1)\Td_j (\dot{\Gamma}, \bp\nabla X,\nabla X + R,\cdots)(\omega+\theta)^{n+1-j} \nonumber\\
&\qquad +j(n+1-j)\Td_j (\dot{\Gamma},\nabla X + R,\cdots)\bp\theta(\omega+\theta)^{n-j} \nonumber\\
&\qquad +(n+1-j)\Td_j(\nabla X + R)\iota_X \a (\omega+\theta)^{n-j} \nonumber\\
&\qquad -(n+1-j)\bp\Td_j(\nabla X + R) \a (\omega+\theta)^{n-j} \nonumber\\
&\qquad +(n+1-j)(n-j)\Td_j(\nabla X + R) \a \bp\theta(\omega+\theta)^{n-1-j},
\end{align}
where to get the last identity we used $\bp R=\bp\omega=0$ and integration by parts.

Now by definition of $\theta$ and Lemma \ref{lulemm}, we have
\begin{equation}\label{lem1}
\iota_X (\omega+\theta)=\iota_X \omega = -\bp\theta
\end{equation}
and
\begin{equation}\label{lem2}
\iota_X (\nabla X+ R)=\iota_X R = -\bp\nabla X.
\end{equation}
Using these, we compute
\begin{align}\label{6}
&\int_M \iota_X \left[j\Td_j(\dot{\Gamma},\nabla X+R,\cdots)(\omega+\theta)^{n+1-j}+(n+1-j)\Td_j(\nabla X+R)\a(\omega+\theta)^{n-j} \right]\nonumber\\
&=\int_M j\Td_j(\iota_X \dot{\Gamma},\nabla X+R,\cdots)(\omega+\theta)^{n+1-j}\nonumber\\
&\qquad-j(j-1)\Td_j(\dot{\Gamma},-\bp\nabla X,\nabla X+R,\cdots)(\omega+\theta)^{n+1-j}\nonumber\\
&\qquad -j(n+1-j)\Td_j (\dot{\Gamma},\nabla X+R,\cdots)(-\bp\theta)(\omega+\theta)^{n-j} \nonumber\\
&\qquad +j(n+1-j)\Td_j (-\bp\nabla X,\nabla X+R,\cdots)\a(\omega+\theta)^{n-j} \nonumber\\
&\qquad +(n+1-j) \Td_j (\nabla X+R) \iota_X \a (\omega+\theta)^{n-j} \nonumber\\
&\qquad -(n+1-j)(n-j)\Td_j (\nabla X+R) \a (-\bp\theta)(\omega+\theta)^{n-1-j}.
\end{align}
Since $\bp \Td_j(\nabla X + R)=j\Td_j(\bp\nabla X,\nabla X + R,\cdots)$, (\ref{5}) and (\ref{6}) are equal.

Hence we obtain for some differential form $\eta$,
\begin{equation*}
\frac{d}{dt}\widetilde{F_j}(\omega_t,X)=\int_M \iota_X \eta =0
\end{equation*}
by dimensional reason.

We now turn to the proof for (\ref{FbarX}). Note that 
\begin{equation*}
\Td_j(R+(\overline{\nabla X})^{\flat,\sharp})=\Td_j(\overline{R}+\overline{\nabla X}),
\end{equation*}
where on the right hand side trace is taken as $\End(T^{0,1}M)$.
Then as the proof for (\ref{FX}) shows, we only need to check the following identities:
\begin{equation*}
\dot{\overline{R}}=-\p \dot{\overline{\Gamma}};\quad \dot{\overline{\nabla X}}=\iota_{\overline{X}} \dot{\overline{\Gamma}};\quad \dot{\omega}=\p \overline{\alpha};\quad \dot{\overline{\theta}}=\iota_{\overline{X}}\overline{\alpha},
\end{equation*}
and
\begin{equation*}
\iota_{\overline{X}} (\omega-\overline{\theta})= -\p\overline{\theta};\quad 
\iota_{\overline{X}} (\overline{R}+\overline{\nabla X})=\p \overline{\nabla X}.
\end{equation*}
For example,
\begin{equation*}
\dot{\overline{R}}=\dot{R}^{\bar{q}}{}_{\bar{p}r\bl}dz^r\wedge d\bar{z}^l=\dot{\overline{R_{p}{}^{q}{}_{l \overline{r} } }}dz^r\wedge d\bar{z}^l=\overline{ -\p_{\br} \dot{\Gamma^{q}_{pl}}}dz^r\wedge d\bar{z}^l=-\p_r \dot{\overline{\Gamma^{q}_{pl}}}dz^r\wedge d\bar{z}^l=-\p\dot{\overline{\Gamma}}.
\end{equation*}
The proof then goes without changes, and one can show that the time derivative is equal to 
\begin{equation*}
\int_M \iota_{\overline{X}} \left[j\Td_j(\dot{\overline{\Gamma}},\overline{\nabla X}+\overline{R},\cdots)(\omega-\overline{\theta})^{n+1-j}-(n+1-j)\Td_j(\overline{\nabla X}+\overline{R})\overline{\a}(\omega-\overline{\theta})^{n-j} \right],
\end{equation*}
which is zero. Note the minus sign in front of the second term.
\end{proof}

\end{document}
